\section{Introdução}

O mercado automotivo global e brasileiro atravessa uma transição tecnológica impulsionada pela entrada acelerada de Veículos Elétricos (EVs) \cite{alvarez2024modelagem}. Diferente dos veículos a combustão interna (ICE), cujo comportamento de depreciação histórica é consolidado e linearizado pela idade do bem, os modelos eletrificados introduzem variáveis de incerteza operacional: degradação acelerada de baterias, rápida obsolescência tecnológica, políticas tarifárias voláteis e guerras de preços \cite{gautam2024depreciation, li2026global}.

Compreender e segmentar essas trajetórias temporais não é apenas uma demanda mercadológica de precificação e análise de risco para seguradoras e instituições financeiras; constitui um problema clássico de \textit{Data Mining} voltado à mineração de séries temporais (\textit{time-series mining}) e descoberta de padrões em dados ordenados. Métodos convencionais frequentemente achatam as curvas históricas ou dependem de distâncias euclidianas 1D que falham em capturar dinâmicas não lineares, flutuações sazonais e correlações defasadas entre segmentos concorrentes \cite{schloter2022empirical}. Embora medidas de similaridade elásticas específicas para esse tipo de dado, como o \textit{Dynamic Time Warping} (DTW), ofereçam maior flexibilidade ao alinhar sequências defasadas temporalmente, elas ainda operam no domínio unidimensional e frequentemente apresentam alto custo computacional em grandes volumes de dados.

A representação convencional de séries temporais como vetores unidimensionais impõe limitações analíticas severas ao processo de mineração de dados, restringindo os algoritmos de agrupamento a métricas de distância altamente suscetíveis a variações de escala. Essa restrição metodológica fundamenta a necessidade de investigar transformações de atributos mais robustas, avaliando se representações bidimensionais geradas por \textit{Gramian Angular Fields} (GAF) ou \textit{Markov Transition Fields} (MTF) são capazes de codificar correlações cruzadas e dependências dinâmicas que escapam às métricas clássicas, como a distância Euclidiana e o DTW \cite{wu2023imaging, ulyanin2019feature}.

Sob o ponto de vista da modelagem estatística e do aprendizado de máquina, o estudo se justifica ao contrapor a eficácia de espaços latentes extraídos dessas imagens frente a métodos tradicionais de redução de dimensionalidade e seleção de atributos. A aplicação de matrizes GAF/MTF permite transpor o problema de dados tabulares temporais para o domínio de operadores de extração de padrões espaciais, abrindo caminho para testar se técnicas de aprendizado não supervisionado se beneficiam da preservação topológica dessas transformações \cite{jueco2024performance}. Com isso, o trabalho estabelece um teste empírico controlado para mensurar se a complexidade adicional dessa projeção espacial se traduz, de fato, em \textit{clusters} mais coesos e discriminantes do que os obtidos via abordagens vetoriais convencionais.

Por fim, o trabalho preenche uma lacuna relevante na análise de fenômenos de mercado ao estruturar a descoberta de padrões de desvalorização e possíveis efeitos de transbordamento decorrentes da chegada dos veículos elétricos e das montadoras chinesas \cite{li2025rise, abrardi2026market}. Ao analisar o co-movimento de preços entre matrizes energéticas concorrentes, a pesquisa vai além do agrupamento meramente exploratório e fornece um protocolo experimental rigoroso baseado em técnicas consolidadas de pré-processamento, mitigação do impacto de \textit{outliers} e validação de modelos. Dessa forma, a investigação entrega simultaneamente uma contribuição teórica sobre representação de dados ordenados na mineração de dados e uma ferramenta descritiva aplicada para precificação e gestão de risco sobre ativos automotivos.
